\section{Interfaces with arbitrary endpoint degrees}\label{sec:generalports}

Here a \emph{root} means a vertex of original degree at least three,
without an upper bound. We use the same decomposition into direct root
edges, ordinary paths, return paths and pendant paths as before. A
non-direct ordinary path has two different actual end edges; their
weights may be chosen independently. A direct edge has one shared weight.
This distinction is essential throughout.

For a path $x_1=a,\ldots,x_\ell=b$ with internal vertices of degree two,
prescribed final endpoint counts $(r,s)$ give the allowed next weights
\begin{equation}\label{eq:generalP}
 P_{r,s}(a)=\{0,1,2\}\setminus
 \bigl(\{r-a:a>0\}\cup\{4-a-s:a<2\}\bigr).
\end{equation}
Values outside $\{0,1,2\}$ have no effect. Every internal comparison is
equivalent to $x_i\ne x_{i+2}$. Thus the odd and even positions form
independent proper three-colour paths with endpoint constraints.

\begin{lemma}\label{lem:generalpath}
If the two allowed sets in~\eqref{eq:generalP} each have at least two
elements, any end weights extend at length at least four; at length
three they extend if and only if $a\ne b$. Opposite pure end weights
extend at every length at least three for all physical endpoint counts.
If all four endpoint colour degrees are at least four, distinct end
weights extend already at length two.
\end{lemma}
\begin{proof}
At length three choose the middle value in the intersection of two
two-element subsets of a three-element set. The other condition is
$a\ne b$. At length four choose $x_2\in P_L\setminus\{b\}$ and
$x_3\in P_R\setminus\{a\}$. At length five choose
$x_3\notin\{a,b\}$ and different $x_2\in P_L,x_4\in P_R$.
For greater lengths each position path has at least two edges between
its endpoint constraints. A three-colour path of at least two edges
connects arbitrary prescribed colours, by induction starting with a
middle colour different from both ends. For a pure end weight, at most
one value is forbidden, proving the opposite-pure assertion. If
$r,s\ge4$, a mixed port forbids no value, and a pure port can only
forbid its own value, when the relevant colour degree is four. At
length two distinct ports therefore satisfy both endpoint checks.
\end{proof}

We will also use the following small extension of the first part: at
length at least three, different end weights extend if one allowed set
is $\{1\}$, the other contains 1 and has at least two elements, and the
opposite fixed end weight is pure. At lengths three and four this
follows from the same explicit choices; at length five the two allowed
even-position values can be made different; greater lengths follow
from the position-path induction.

\subsection{Exact preservation under subdivision}\label{sec:preserve}

\begin{theorem}\label{thm:preserve}
Let a legal colouring of $\doub Q$ be given. Replace a set of original
edges $e=uv$ by internally disjoint paths of lengths $L_e\ge4$. Set
\[
 \epsilon_e(v)=\mathbb{1}_{\{x_e=1,\ R_Q(v)=B_Q(v)\in\{1,3\}\}}.
\]
If $L_e\ge4+\epsilon_e(u)+\epsilon_e(v)$ for every replaced edge, the
colouring extends preserving every old vertex colour degree and every
unreplaced edge weight. At lengths four and five this condition is
also necessary when both new end weights must equal $x_e$.
\end{theorem}
\begin{proof}
Keep each end weight equal to $a=x_e$. Pure ports have at least two
allowed next values. For $a=1$, the forbidden next values are $R-1$
and $3-B$. If both belong to $\{0,1,2\}$, then $R,B\in[1,3]$.
Since $R+B$ is even, the allowed set has size one exactly when
$(R,B)=(1,1)$ or $(3,3)$, and that set is $\{1\}$. At length four
each next value must differ from the opposite fixed end value $a$;
one singleton is therefore fatal. At length five the two next values
must differ from one another, so precisely two singletons are fatal.
At length at least six each position path can connect arbitrary
nonempty boundary sets. These conditions give the stated threshold.
All old incident contributions remain identical, so all unreplaced
comparisons persist simultaneously.
\end{proof}

For cubic initial data there are 60 ordered legal end states: for each
$a$, each endpoint red degree lies in $\{a,\ldots,a+4\}$ and the two
degrees differ. All 40 pure-port states extend at length four. Of the
20 mixed-port states, exactly eight fail: one end red degree is three
and the other is one of $1,2,4,5$, in either order. These are fixed-data
obstructions, not counterexamples to Theorem~\ref{thm:main}.

In particular selected edges whose endpoints all have degree at least
four can always be replaced by paths of length at least four while
preserving a given colouring. Known regular-graph colourability
\cite{GOW26} makes this unconditional for every connected regular
mother graph of degree at least four. Known bipartite colourability
\cite{GW23} gives the same conclusion for a bipartite mother graph
when every selected endpoint has degree at least four. The initial
colouring is supplied by those theorems; it is not assumed for a
general unknown mother graph.

If a vertex-minimal counterexample to the general conjecture exists,
it has no ordinary degree-two thread of length at least five between
distinct nonadjacent roots; if both roots have degree at least four,
length at least four is excluded. Replace that thread by a single edge.
The smaller graph is simple and is not $K_2$. Its legal colouring
exists by minimality. On a mixed new edge at most one endpoint can have
$(R,B)=(3,3)$, since two such endpoints would already violate the edge.
Theorem~\ref{thm:preserve} restores the thread, a contradiction. The
nonadjacency assumption cannot be dropped in this argument.

\section{One-bridge trees and decorated paths}\label{sec:trees}

For a rooted tree $(T,r)$ and an external bridge of weight $c$, let
$P_c(T,r)$ be the achievable \emph{internal} root red sums when all
internal edges are legal after the root receives external contribution
$(c,2-c)$. The external bridge itself is not yet checked.

\begin{lemma}\label{lem:treeports}
For a one-vertex tree all three sets are $\{0\}$. For a two-vertex tree,
$P_0=\{0\},P_1=\{0,1,2\},P_2=\{2\}$. Every tree of order at least two
has $|P_1|\ge3$; every tree of order at least three has
$|P_0|,|P_2|\ge2$.
\end{lemma}
\begin{proof}
Use simultaneous induction on tree order. Delete the root, leaving
$k$ rooted subtrees with root internal degrees $d_i$. If the root edge
to subtree $i$ has weight $a_i$ and its internal root sum is $p$, the
only new restrictions, for final parent counts $A,B$, are
\[
 a_i>0\Rightarrow p+a_i\ne A,\qquad
 a_i<2\Rightarrow2d_i-p+2-a_i\ne B.
\]
Thus a mixed edge forbids at most two child candidates, a pure edge at
most one. Colour complementation gives
$P_{2-c}=\{2k-p:p\in P_c\}$.

For $c=0$, put all $a_i=1$. The internal sum $k$ is feasible: larger
children have three mixed candidates, and singleton children have
counts $(1,1)$, different from parent counts $(k,k+2)$ unless $k=1$,
which would be the two-vertex base case. Change one $a_i$ to 2 to seek
sum $k+1$ and parent counts $(k+1,k+1)$. That child must avoid
$p=k-1$ in $P_2$. A child of order at least three has two candidates;
a singleton fails only at the excluded $k=1$; a two-vertex child fails
only at $k=3$. The only remaining configuration is three two-vertex
children. Take root weights $(0,1,1)$ and all child internal weights
zero instead. The parent is $(2,6)$; the pure-blue child root has blue
degree four and the mixed child roots have counts $(1,3)$. This gives
a second sum distinct from $k=3$. Complementation proves the $P_2$ bound.

For $c=1$, all $a_i=1$ yield sum $k$ and counts $(k+1,k+1)$.
Changing one $a_i$ to zero seeks sum $k-1$ and counts $(k,k+2)$;
the child must avoid $p=2d_i-k$ in $P_0$. The only exception not
handled by the induction is two two-vertex children. In that case set
both root weights and both child internal weights to zero, giving
internal sum zero and parent counts $(1,5)$, with child blue degrees
four. Consequently a feasible sum $a\ne k$ always exists. Its colour
complement $2k-a$, together with $k$, gives three distinct sums.
\end{proof}

For an outside vertex of final counts $(A,B)$ the bridge imposes at
most one or two further exclusions. Lemma~\ref{lem:treeports} gives the
complete compatibility table
\begin{center}
\begin{tabular}{lccc}\toprule
Tree order&$c=0$&$c=1$&$c=2$\\\midrule
1&$B\ne2$&$A,B\ne1$&$A\ne2$\\
2&$B\ne4$&always&$A\ne4$\\
at least 3&always&always&always\\\bottomrule
\end{tabular}
\end{center}
Here ``always'' assumes physical counts containing the bridge copies.
This is an interface statement; adding the bridge changes outside
degrees and may invalidate other outside edges.

\subsection{Saturation of decorated paths}\label{sec:decorated}

\begin{theorem}\label{thm:decorated}
Let every internal vertex of a path have at most one bridge to a
disjoint rooted tree. A path of at least six edges accepts arbitrary
physical final endpoint counts and arbitrary prescribed first and last
edge weights, while colouring all attached trees. The endpoints may be
the same vertex if both incident contributions are counted.
\end{theorem}
\begin{proof}
Classify an internal vertex as $A$ (bare), $L$ (a singleton tree),
$S$ (a two-vertex tree), or $F$ (a larger tree). With incoming and
outgoing weights $a,p$, bridge weight $c$, and degree $D=2$ or 3, its
red sum is $r=a+p+c$. For $A$, put $c=0$. The compatibility table
above says: for $L$, $a+p=0$ forces $c=0$, $a+p=4$ forces $c=2$,
and otherwise any $c$ works; for $S$, $a+p=2$ forces $c=1$ and
otherwise any $c$ works; for $F$ every $c$ works.

Record a state $(p,D,r)$. There are 24 possible states,
\[
 \mathcal U=\{(p,2,r):p\le r\le p+2\}
 \cup\{(p,3,r):p\le r\le p+4\},\quad p\in\{0,1,2\}.
\]
An old state with outgoing weight $a$ can precede a new state
$(p,D',r')$ exactly when the indicated bridge is allowed and
$a>0\Rightarrow r\ne r'$, $a<2\Rightarrow2D-r\ne2D'-r'$.
Let $T_X$ be the resulting monotone map on state subsets. Complete
integer tables are given in Appendix~\ref{app:decorated}, including the
four lower-bound families $F_1,F_2,F_3,F_4$ of sizes $30,19,7,3$.

The initial endpoint has just 93 comparison types: each colour count
is $0,\ldots,6$ or greater than six, with physical and parity
restrictions. For each type and first letter the reachable set contains
an element of $F_1$. For every $M\in F_i$ and letter $X$, $T_X(M)$
contains an element of $F_{i+1}$, for $i=1,2,3$.
For each of the three elements of $F_4$, applying $T_X$ gives the
entire locally feasible state set $\Omega_X$. These are complete finite
identities, recomputed without the constructor by
\code{decorated_path_certificate_check.py}.

It follows by monotonicity and induction that after five internal
vertices the reachable set is $\Omega_X$, and remains saturated after
every further vertex. For a prescribed final weight, $\Omega_A$ has
three possible red sums, $\Omega_L,\Omega_F$ have five and $\Omega_S$
has at least four. The right boundary forbids at most two, so it can
always be met. Restore the attached trees using the bridge lemma.
The clipping of large endpoint counts is exact, since internal colour
degrees never exceed six. This proves arbitrary length and arbitrary
tree shapes from the finite identities.
\end{proof}

\section{Sparse direct branch cores}\label{sec:branches}

\begin{theorem}[A2]\label{thm:branches}
Let $W=\{v:d_G(v)\ge3\}$ for a connected simple graph $G\ne K_2$.
Each of the following hypotheses implies $\lir(\doub G)\le2$:
\begin{enumerate}
\item $G[W]$ is independent;
\item there is no degree-three vertex and $\Delta(G[W])\le1$;
\item $\Delta(G[W])\le1$ and every non-direct ordinary branch path
has length at least four.
\end{enumerate}
Return paths of length at least three and pendant paths of arbitrary
positive length are allowed in all three statements.
\end{theorem}

\subsection{Independent branch vertices}
\begin{proof}[Proof of the first statement]
For a root let $l,h,p$ count return paths, ordinary incidences and
pendant paths, respectively; its original degree is $D=2l+h+p$.
On $W$ form a simple graph joining distinct roots connected by an
ordinary path of length at least three. Choose a maximal independent
set $A$, then a maximal independent set $B$ in the remaining graph,
and put $C=W\setminus(A\cup B)$. Every $C$ root has a long neighbour
in both $A$ and $B$; every isolated root belongs to $A$.

Give every return path ports $(0,2)$. Give pendant root ports
0, 2, 1 on $A,B,C$, respectively. On ordinary paths assign port pairs
$AA:(0,0)$, $BB:(2,2)$, $AB:(0,2)$, $AC:(0,1)$, $BC:(2,1)$.
An $AA$ or $BB$ path necessarily has length two. Balance-orient the
multigraph of $CC$ paths and give each tail port 0 and head port 2.
For a length-two $AC$ path whose $C$ root has degree three, change
its $C$ port to 2; for such a $BC$ path change it to zero.

At a $C$ root of degree at least four the only imbalance between
its red sum and $D$ is at most one from the balanced $CC$ orientation.
At degree three the two long $A,B$ incidences leave at most one other
incidence, so either one $CC$ imbalance or one short-path correction
occurs, again of magnitude at most one. If the red sum is $D-1$,
change its long $A$ port from 1 to 2. If it is $D+1$, change its
long $B$ port from 1 to zero. These independent end ports affect no
other root. Now every $C$ root has $(R,B)=(D,D)$.

At an $A$ root the counts are $(2l,2(D-l))$, and at a $B$ root they
are reversed. The larger count is an even integer at least four.
Equality to four at a root with an external path occurs only when
$D=3,l=1,h+p=1$. Such a root is exceptional only in the following
two cases. If its unique external path has length two and joins
another $A$ root, change its own port from 0 to 2; if both ends
are exceptional change just one. Its counts become $(4,2)$ and the
midpoint has $(2,2)$, legal on the red and blue sides. If its unique
external path is a two-edge pendant, change that port to 1 and use
the word $(1,1)$; root, middle and leaf counts are $(3,3),(2,2),(1,1)$.
These exceptional roots have no long ordinary incidence, so none of
the previous adjustments is disturbed.

The remaining two-edge ordinary checks are listed below. The displayed
root counts refer to the colour present on the corresponding port.
\begin{center}\small
\begin{tabular}{llll}\toprule
Type&Ports&Middle $(R,B)$&Root comparisons\\\midrule
$AA$&$0,0$&$0,4$&both blue at least 6\\
$BB$&$2,2$&$4,0$&both red at least 6\\
$AB$&$0,2$&$2,2$&relevant counts at least 4\\
$AC$, $D_C=3$&$0,2$&$2,2$&blue at least 4; red 3\\
$AC$, $D_C\ge4$&$0,1$&$1,3$&blue at least 4; both at least 4\\
$BC$&complement&complement&same inequalities\\
$CC$&$0,2$ or $2,0$&$2,2$&both root counts at least 3\\\bottomrule
\end{tabular}
\end{center}
For long ordinary paths the end ports differ. An $A$ or $B$ allowed
set has at least two elements and contains 1. At a $C$ root the only
singleton occurs at $D=3$ with port 1, when it is $\{1\}$; its
other end is an $A$ or $B$ pure port. The general path lemma and its
singleton extension therefore apply. Every return path has opposite
pure ports. Pendant pure ports use the constant word at lengths one
or two, with the unique bad two-edge case already repaired; at longer
lengths use the opposite pure leaf port. A $C$ pendant uses $(1)$ or
$(1,1)$ at lengths one or two; otherwise give its leaf port zero and
apply the singleton extension. All root counts were fixed first.
Consequently the independent path fillings are simultaneously legal.
\end{proof}

\subsection{A direct matching with no degree-three roots}
\begin{proof}[Proof of the second statement]
Compress all paths, using a private placeholder leaf for each pendant
path only for the purpose of orientation. A balanced orientation of
this multigraph gives outdegrees $o(v)\in\{\lfloor D/2\rfloor,
\lceil D/2\rceil\}$. Give non-direct tail ports 1 and head ports 0;
give a pendant root port 1 or 0 according to its direction; give each
direct edge the shared weight 1. If $i_H(v)$ indicates an incoming
direct edge, the initial counts are
\[
 R_0(v)=o(v)+i_H(v)\ge2,\qquad
 B_0(v)=2D-R_0(v)\ge5.
\]
Repair each direct matching edge independently. If only its red
degrees coincide, change its weight to zero. The head initially had
red degree at least three, so both new red degrees remain at least
two. If only its blue degrees coincide, change its weight to 2.
The tail initially had blue degree at least six, so both new blue
degrees remain at least five. If neither coincides, retain weight 1.

If both coincide, orient the edge $u\to v$. Equal original degrees
and $o(u)=o(v)+1$ force $D=2k+1$, $k\ge2$, with common initial
counts $(k+1,3k+1)$. There are $k$ other outgoing incidences at $u$.
If one is an ordinary or return path, change its tail port 1 to 2,
giving $u$ counts $(k+2,3k)$. Otherwise choose an outgoing pendant
port and change 1 to zero, giving $(k,3k+2)$. The mixed direct edge
is now legal. Only tail ports were modified; head ports stay zero,
so even repairs selecting the same connection cannot interfere.
Different direct edges have disjoint endpoint sets.

Final root counts satisfy $R\ge2,B\ge5$, and a non-direct port 2
occurs only at $R\ge4$. Every non-direct port pair is $(0,1)$ or
$(0,2)$ up to reversal. Its two allowed sets have size at least two,
so every length at least three extends. At length two the midpoint
is respectively $(1,3)$ or $(2,2)$ and the same inequalities verify
both ends. This includes return paths at their legal lengths.
Pendant port zero uses $(0)$, $(0,0)$, or opposite leaf port 2;
pendant port one uses $(1)$, $(1,0)$, or leaf port zero, according
as its length is one, two, or at least three. Direct checking or
Lemma~\ref{lem:generalpath} verifies every case.
\end{proof}

\subsection{A direct matching with ordinary paths of length four}
\begin{proof}[Proof of the third statement]
Let $h\in\{0,1\}$ indicate a direct incidence, let $l$ count return
paths, and let $z$ count all other root incidences. Thus $D=h+2l+z$.
Give direct edges weight 1, return ports $(0,2)$, and choose the
$z$ free ports from $\{0,2\}$. For a root with $h=1$, if $a$ free
ports are 2 its counts are
\[
 (R,B)=(1+2l+2a,\ 1+2l+2(z-a)).
\]
For a direct edge, if both ends have free ports and one has $z\ge2$,
fix the other end and try $a=0,1,2$. At most two candidates are
forbidden. If both have $z=1$, their candidates are
$(D-1,D+1)$ and $(D+1,D-1)$: use the same choice at different degrees
and opposite choices at equal degrees. If one end has $z=0$, it has
a return path. Change one such port pair from $(0,2)$ to $(0,1)$.
That end now has even counts $(D-1,D+1)$ and the other keeps odd
counts. If both have $z=0$, change just one end. Thus every direct
edge is legal, independently of all others.

At roots with $h=0$, leave free ports zero; counts are
$(2l,2(l+z))$. The sole exceptional pendant is $l=z=1$ with length
two. Change its port to 1, making the root $(3,3)$, and use $(1,1)$.
Every ordinary non-direct path has pure ports and length at least
four, so the general path lemma applies even when its ports are equal.
Unmodified returns have ports $(0,2)$. A modified return has ports
$(0,1)$ at counts $(D-1,D+1)$; both allowed sets have at least two
elements, so it also extends at length three. At a directly matched
root, each colour used by a pendant pure port has odd degree at least
three, different from the short pendant values two and four. At an
unmatched root its blue degree is at least four, and equality to four
on a two-edge pendant was exactly the exception repaired above.
Long pendants use opposite pure leaf ports. These cases exhaust the
construction.
\end{proof}

In all three proofs the case $W=\varnothing$ is separate: a single
vertex has no edges; a two-edge path has word $(0,0)$; a longer path
uses opposite pure leaf ports. For a cycle, properly three-colour the
position graph joining positions $i$ and $i+2$ cyclically. Its maximum
degree is at most two and its resulting word satisfies every internal
comparison. The excluded one-edge path is $K_2$.

\section{Homogeneous direct cores with long connections}\label{sec:homogeneous}

\begin{lemma}[Even signed forms]\label{lem:evenforms}
Suppose each of finitely many homogeneous linear forms has an even
number, at least four, of distinct variables, all coefficients are
$\pm1$, and each variable occurs in at most two forms. There is a
choice of signs $\pm1$ for the variables making every form nonzero.
\end{lemma}
\begin{proof}
Make a signed multigraph whose vertices are forms and whose shared
variables are edges; variables occurring once are boundary variables.
In a component with a boundary variable, retain a spanning tree and
one root boundary variable, set all other variables to 1, and process
leaves upwards. The one remaining sign at each step avoids at most
one forbidden value.

In a component without boundary variables every vertex has even
degree at least four. Give each edge the product of its endpoint
coefficients as sign. There is a positive cycle: if no fundamental
cycle is positive, two negative fundamental cycles sharing a tree
edge have a positive cycle as their symmetric difference; if none
share tree edges the graph is a cactus and has a vertex of degree at
most two, a contradiction. Keep such a cycle and a forest rooted on
it. Eliminate off-cycle vertices as before. On the cycle the relations
are $a_i y_i+b_i y_{i-1}+C_i\ne0$ with even $C_i$.
For $C_i=0$ the allowed relation is a bijection on two signs; for
$C_i=\pm2$ it has three of the four pairs; for $|C_i|\ge4$ all four
pairs. If every relation is a bijection, positivity makes their
composition have a fixed point. Otherwise their composition has at
least three pairs: bijections preserve size, and composition of
relations with at least three pairs has at least three pairs. Such
a relation contains a diagonal pair. Hence the signs close around
the cycle.
\end{proof}

\begin{theorem}[A3]\label{thm:homogeneous}
Let $G$ have minimum degree at least two. In each component of
$H=G[\{v:d_G(v)\ge3\}]$, suppose the original degree is constant and
the component is three-colourable. Suppose all ordinary non-direct
connections have length at least three and there are no return paths.
Then $\lir(\doub G)\le2$.
\end{theorem}
\begin{proof}
Orient each long connection and give it opposite pure ports. Its
contribution $c(v)$ is twice an indegree. For an isolated direct-edge
component $uv$ of original degree $D$, the two long-incidence counts
are both $D-1$. Reversals of the long paths are sign variables, and
$c(u)-c(v)$ is an even form with $2(D-1)\ge4$ incidences. If a long
path joins $u$ and $v$ its sign occurs twice with combined coefficient
$\pm2$; it is private to this equation. Temporarily omit equations
with such variables. All remaining forms satisfy
Lemma~\ref{lem:evenforms}; then restore each omitted equation by its
private variable, one of whose two signs makes it nonzero. Give each
isolated direct edge weight 1. It is now legal in both colours.

For every larger direct component choose proper targets $\phi$ over
$\mathbb F_3$ and solve the internal incidence equations
$\sum_{e\ni v}x_e=\phi(v)-c(v)$. A nonbipartite component realises
arbitrary demands: a spanning tree and an edge joining two vertices
of the same tree parity allow adjustment of the signed total by an
invertible coefficient $\pm2$, followed by leaf elimination. In a
bipartite component solvability is precisely the alternating-sum
condition in Lemma~\ref{lem:linear}. Choose the larger side, of size
at least two, initially all zero and the other side all one. Changing
zero, one or two vertices of the first side to 2 changes that sum by
$0,2,4$, covering all residues. The targets stay proper. Representatives
$0,1,2$ of the solved edge weights give different red sums on every
direct edge. The original degrees agree there, so the blue sums also
differ. Isolated roots require no direct comparison. Finally fill
all long paths with their opposite pure ports.
\end{proof}

For an $r$-regular simple mother graph, replacing selected edges by
paths of length at least three is therefore allowed whenever the
retained graph is three-colourable. At $r=3$, a nonempty selected set
leaves every retained component 2-degenerate, since a cubic retained
component would be a whole connected component of the original graph.
This recovers the earlier cubic subdivision constructions, already
contained in A1. At unbounded degree the three-colourability hypothesis
is substantive and has not been removed.

\section{Balanced half-edges and simultaneous extensions}\label{sec:balanced}

\begin{lemma}\label{lem:halfedge}
Let $J$ be a multigraph, with loops allowed and no vertex of degree
one. Its incidences admit values $0,1,2$ summing to $d_J(v)$ at every
vertex, different at the two ends of every edge. Odd-degree vertices
have exactly one incidence of value 1 and even-degree vertices none.
For disjoint pairs of positive-degree vertices of the same parity,
one may instead require exactly one vertex of each pair to have sum
$d_J(v)\pm1$, leaving every other sum unchanged.
\end{lemma}
\begin{proof}
Balance-orient $J$, assigning 0 at tails and 2 at heads. At an odd
vertex choose one majority-valued incidence to change to 1. There
are $\lceil d/2\rceil\ge2$ different candidate edge identifiers;
a loop cannot supply two majority-valued candidates. The request-edge
bipartite graph has each request adjacent to at least two edges, each
edge to at most two requests. Hall's condition follows by counting,
so all changes can use distinct edges. The sums become exact and
end values stay different.

For an odd-degree pair skip the correction at one member. For an
even-degree pair create one request selecting any incidence at either
member and changing it to 1; the union contains at least two edge
identifiers. Each original vertex belongs to at most one request,
so the same Hall argument applies. These changes give precisely the
claimed one-unit deviations.
\end{proof}

\begin{theorem}[A4]\label{thm:partialbalanced}
Let a connected multigraph $M$ have original degrees $D(v)\ge4$.
Subdivide an edge set $J$ into paths of length at least two, with all
loops in $J$ and loop paths of length at least three. Suppose $J$ has
no vertex of degree one, the retained graph $H=M-J$ is simple, and
every retained edge has endpoints of different original degrees.
Then the simple expanded graph has a legal doubled two-colouring
with $R(v)=B(v)=D(v)$ at every old vertex.
\end{theorem}
\begin{proof}
Use Lemma~\ref{lem:halfedge} on $J$ and give every retained edge
weight 1. Each old vertex receives $d_J+d_H=D$ in both colours.
Retained edges compare different $D$ values. Each subdivided path
has distinct end weights and all four endpoint colour degrees at
least four, so Lemma~\ref{lem:generalpath} fills it at every allowed
length without changing the old sums.
\end{proof}

\begin{theorem}[A5]\label{thm:simultaneous}
Let $H$ be a simple graph whose components of order at least three
have given legal doubled two-colourings. On the same old vertex set
let $J$ be a connection multigraph with $k(v)=d_J(v)$ either zero
or at least four, constant on each component of $H$. Require $k\ge4$
on every $K_2$ component. Replace $J$ edges by internally disjoint
paths of length at least two, or at least three for loops. The
colouring extends to the simple expanded graph preserving every
given old edge weight. No initial colouring of $K_2$ is assumed.
\end{theorem}
\begin{proof}
Pair the endpoints of the $K_2$ components in
Lemma~\ref{lem:halfedge}. All other old vertices receive exactly
$(k,k)$ from new ports. Every given old colour-degree difference is
unchanged, since $k$ is equal at its two ends. Give each old $K_2$
edge weight 1. Its endpoints then have counts
$(k+1+\epsilon,k+1-\epsilon)$ and $(k+1,k+1)$ for
$\epsilon\in\{-1,1\}$, so that edge is legal. All vertices incident
with new paths have both colour degrees at least four. Different
new end ports can therefore be filled at length two by
Lemma~\ref{lem:generalpath}. All old and new checks hold together.
\end{proof}

A vertex-minimal counterexample cannot admit the latter decomposition
with $J\ne\varnothing$: deleting its new internal vertices gives
smaller old components, while the $K_2$ and isolated cases are already
handled. This does not assert that every graph has such a decomposition.

\section{Near-centre parity assignments}\label{sec:parity}

\begin{lemma}\label{lem:nearparity}
A parity prescription $p:V(G)\to\{0,1\}$ can be realised by weights
$0,1,2$ with $R(v)\equiv B(v)\equiv p(v)\pmod2$ and
$|R(v)-d_G(v)|\le1$ if and only if $\sum_{v\in C}p(v)$ is even
in every component $C$. The edges changed from weight 1 may be
restricted to a forest.
\end{lemma}
\begin{proof}
Necessity is the red handshake identity. Start with all weights 1.
In each spanning tree choose a $T$-join for
$T=\{v:p(v)\ne d_G(v)\pmod2\}$ by leaf elimination; $|T|$ is even.
Pair incidences of the chosen forest at each vertex, leaving one
unpaired at each odd-degree vertex. Its edges decompose into open
trails with those unpaired ends; there is no closed trail in a forest.
Along each trail change weights alternately to 2 and zero. Contributions
cancel at every paired passage, leaving a change of $\pm1$ exactly
at the vertices of $T$. All claims follow, including the exact centred
assignment when $T$ is empty.
\end{proof}

\begin{theorem}[A6]\label{thm:gapparity}
Let $F$ consist of the original edges $uv$ with
$|d_G(u)-d_G(v)|\le2$. If $F$ is bipartite and its bipartition
phases can be chosen to give an even total number of odd labels,
then $\doub G$ has a legal two-colouring. Both colours distinguish
the ends of every original edge.
\end{theorem}
\begin{proof}
Realise those labels by Lemma~\ref{lem:nearparity}. On $F$, both
colour degrees have different parities. Off $F$, the intervals
$[d_G(v)-1,d_G(v)+1]$ at the two ends are disjoint, and contain
both colour degrees. This verifies every comparison.
\end{proof}

The phase condition is exact: an odd-order component of $F$ allows
the total parity to be changed by a flip. If all components have even
order, it holds exactly when an even number of components have two
odd-sized sides. For example, connect the roots of three disjoint
$K_{r_i,r_i}$ modules in a triangle, where the $r_i\ge4$ are even
and pairwise differ by at least four. The graph has no degree-two
threads and is nonbipartite. Its $F$ consists precisely of the three
modules with even sides, so the theorem applies. This condition does
not cover general nonbipartite regular cores.
